\begin{titlingpage}
    \begin{center}
        \vspace*{1cm}

        \Huge
        \textbf{Syntaxes with Binding, \\ Their Programs, and Proofs}

        \vspace{1.5cm}

        \LARGE

        \textbf{Guillaume Xavier ALLAIS}

        \vfill

        A thesis presented for the degree of\\
        Doctor of Philosophy

        \vspace{0.8cm}

        \includegraphics[width=0.4\textwidth]{strath_coat}

        \vspace{0.8cm}

        \Large
        Computer and Information Sciences\\
        University of Strathclyde\\
        United Kingdom\\
        2022

    \end{center}
\end{titlingpage}


\pagebreak
{
\hspace{0pt}
\vfill{}
This thesis is the result of the author’s original research. It has been
composed by the author and has not been previously submitted for
examination which has led to the award of a degree.

The copyright of this thesis belongs to the author under the terms of the
United Kingdom Copyright Acts as qualified by University of Strathclyde
Regulation 3.50. Due acknowledgement must always be made of the use of
any material contained in, or derived from, this thesis.

\vspace{1.5cm}

\hfill
\begin{minipage}{0.4\textwidth}
Signed: Guillaume ALLAIS \\
Date: \today{}
\end{minipage}
\vfill
\hspace{0pt}
}
\pagebreak

\section*{Acknowledgements}

I would like to thank my supervisor Conor McBride for his contagious
enthusiasm, for his generosity, and most of all for the near limitless
freedom I was given to pursue the problems I was interested in.
%
I learned a lot of the techniques used in this thesis through our
long discussions.
%
Many thanks naturally go to my other co-authors.
%
James Chapman, James McKinna, and Bob Atkey helped me simplify my
generic setup and clarify some of the categorical structures
underlying this work.
%
Andreas Abel, Steven Schäfer, Kathrin Stark, Aliya Hameer,
Alberto Momigliano, and Brigitte Pientka offered,
via the POPLMark Reloaded Challenge, a great opportunity to
compare, contrast, and mutually enrich our respective approaches
to well scoped syntaxes with binding.

I would like to thank my examiners Neel Krishnaswami and Jules Hedges
for the detailed feedback that helped make this thesis a better document,
and for a very pleasant and interesting discussion during the viva.

The results in this thesis could not have been obtained without the
tireless work of the research team behind Agda (Ulf Norell,
Andreas Abel, Nils Anders Danielsson, Jesper Cockx).
%
I have learned so much from using the system, and have had a fantastic
time contributing back to it and collaborating on new features during
the regular Agda Implementors' Meetings.

My work benefited immensely from Matthew Daggit's development and
maintenance work for the Agda standard library.
Our constant discussions on library issues
have deeply informed my own design sensibilities.

This thesis was a long time in the making.
%
I am immensely grateful to both Bart Jacobs and Edwin Brady for trusting
me to conduct exciting postdoctoral research under their supervision
while I had not yet defended.
%
This also means that I have made many wonderful friends in my time in
Strathclyde, Nijmegen, and St Andrews and I want to conclude
by thanking them for making these years very pleasant.

Last but not least, I want to thank my parents who always supported my
choices.

\pagebreak

\hspace{0pt}
\vfill{}
\begin{abstract}
  Almost every programming language's syntax includes a notion of binder and
  corresponding bound occurrences, along with the accompanying notions of
  α-equivalence, capture-avoiding substitution, typing contexts,
  runtime environments, and so on.

  In the past, implementing and reasoning about programming languages required
  careful handling to maintain the correct behaviour of bound variables. Modern
  programming languages include features that enable constraints like scope
  safety to be expressed in types.
  %
  Nevertheless, the programmer is still forced to write the same boilerplate
  over again for each new implementation of a scope safe operation (e.g.,
  renaming, substitution, desugaring, printing, etc.), and then again
  for correctness proofs.

  In a case study focusing on the simply typed lambda calculus, we analyse these
  well scoped traversals and observe that they all share the same structure.
  %
  This enables us to formulate them as instances of a more general program whose
  properties can be established generically.
  %
  Alas, the programmer is still forced to duplicate this effort for every new
  language they implement.

  This leads us to defining an expressive universe of syntaxes with binding
  and to demonstrate how to implement scope safe traversals once and for all by
  generic programming and how to derive properties of these traversals by
  generic proving.
  %
  Our universe description, generic traversals and proofs, and our examples have
  all been formalised in Agda.
\end{abstract}
\vfill
\hspace{0pt}
\pagebreak
